\loai{Hạt chuyển động vuông góc với đường cảm ứng từ}
\begin{vidu} %[2P4-7.2B][Loại1]
Một điện tích có độ lớn $10\ \mu C$ bay với vận tốc $10^5$ m/s vuông góc với các đường sức vào một từ trường đều có độ lớn cảm ứng từ bằng 1 T. Độ lớn lực Lorentz tác dụng lên điện tích là
\choice
{\True 1 N}
{104 N}
{0,1 N}
{0 N}
\loigiai{Độ lớn của lực Lorentz:
\begin{align*}
f=\left|{q_0}\right|vB\sin \alpha=10.10^{-6}.10^5.1.\sin 90^\circ=1\ N.
\end{align*}}
\end{vidu}
\loai{Hạt chuyển động xiên góc với đường cảm ứng từ}
\begin{vidu} %[2P4-7.2B][Loại2]
	Một proton bay vào trong từ trường đều theo phương hợp với đường sức $30\degree$ với vận tốc ban đầu ${3.10^7}$ m/s, từ trường $B = 1,5$ T. Lực Lorenxơ tác dụng lên hạt đó là
	\choice
	{\True $3,6.10^{-12}$ N}
	{$1,8.\sqrt3 .10^{-12}$ N}
	{$0,36.10^{-12}$ N}
	{$36.10^{-12}$ N}
	\loigiai{$f = |q|vB\sin \theta = 1,6.10^{-19}.3.10^7.1,5.\sin 30^\circ = 3,6.10^{-12}$ N}
\end{vidu}
\begin{vidu} %[2P4-7.2B][Loại2]
	Một electron chuyển động với vận tốc $2.10^6$ m/s vào trong từ trường đều $B = 0,01\ T$ chịu tác dụng của lực Lorentz $16.10^{-16}$ N. Góc hợp bởi vector vận tốc và hướng đường sức từ trường là
	\choice
	{$60^\circ$}
	{\True $30^\circ$}
	{$90^\circ$}
	{$45^\circ$}
	\loigiai{
		Độ lớn của lực Lorentz:
	\begin{align*}
			f=\left|q_0\right|vB\sin \alpha\Rightarrow \sin \alpha=\dfrac{f}{\left|q_0\right|vB}=\dfrac{16.10^{-16}}{\left|-1,6.10^{-19}\right|.2.10^6.0,01}=0,5\Rightarrow \alpha=30^\circ.
	\end{align*}
}
\end{vidu}
\loai{Mối quan hệ giữa lực và vận tốc của hạt}
\begin{vidu} %[2P4-7.2B][Loại3]
	Một điện tích bay vào một từ trường đều với vận tốc $2.10^5$ m/s thì chịu một lực Lorentz có độ lớn là 10 mN. Nếu điện tích đó giữ nguyên hướng và bay với vận tốc $5.10^5$ m/s vào thì độ lớn lực Lorentz tác dụng lên điện tích là
	\choice
	{\True 25 mN}
	{4 mN}
	{5 mN}
	{10 mN}
	\loigiai{Độ lớn của lực Lorentz:
		\begin{align*}
			f=\left|{q_0}\right|vB\sin \alpha\Rightarrow \dfrac{f_2}{f_1}=\dfrac{v_2}{v_1}\Rightarrow f_2=f_1\dfrac{v_2}{v_1}=25\ mN.
	\end{align*}}
\end{vidu}
\loai{Mối quan hệ giữa lực và độ lớn điện tích của hạt}
\begin{vidu} %[2P4-7.2B][Loại4]
	Hai điện tích $q_1 = 10\ \mu C$ và điện tích $q_2$ bay cùng hướng, cùng vận tốc vào một từ trường đều. Lực Lorentz tác dụng lần lượt lên $q_1$ và $q_2$ là $2.10^{-8}$ N và $5.10^{-8}$ N. Độ lớn của điện tích $q_2$ là
	\choice
	{\True $25\ \mu C$}
	{$2,5\ \mu C$}
	{$4\ \mu C$}
	{$10\ \mu C$}
	\loigiai{Độ lớn của lực Lorentz:
		\begin{align*}
			f=\left|{q_0}\right|vB\sin \alpha\Rightarrow \dfrac{f_2}{f_1}=\dfrac{\left|{q_2}\right|}{\left|{q_1}\right|}\Rightarrow \left|{q_2}\right|=\left|{q_1}\right|\dfrac{f_2}{f_1}=25\ \mu C.
	\end{align*}}
\end{vidu}
